\centerline{\bf \Huge Итоговая Олимпиада}

\begin{flushright}
        \scriptsize{Врождённый талант может не принести счастья. Счастливы те, кто верят в себя до самого конца и упорно трудятся.}
\end{flushright}

\problem{\textbf{1}} Сколько существует пар натуральных чисел, у которых наименьшее общее кратное (НОК) равно 2000?


\problem{\textbf{2}} Из натурального числа вычли сумму его цифр и получили 2007. Каким могло быть исходное число? Перечислите все варианты!

\problem{\textbf{3}}  В шеренге стоят 2025 человек, среди которых один Начальник. Каждый из стоящих в шеренге либо рыцарь, который всегда говорит правду, либо лжец, который всегда лжёт. Каждый, кроме Начальника, сказал: "Между мной и Начальником стоят ровно два лжеца". Сколько лжецов в этой шеренге, если известно, что Начальник – рыцарь?

\problem{\textbf{4}} Можно ли вычеркнуть одно из натуральных чисел от 1 до 9, а
оставшиеся числа расставить в вершинах куба так, чтобы суммы чисел на каждой грани куба
были равны между собой, но не были кратны вычеркнутому числу?

\problem{\textbf{5}} В пять горшочков, стоящих в ряд, Кролик налил три килограмма мёда (не обязательно в каждый и не обязательно поровну). Винни-Пух может взять любые два горшочка, стоящие рядом. Какое наибольшее количество мёда сможет \textbf{гарантированно} (вне зависимости от того, как налил мёд Кролик) съесть Винни-Пух?
